% Standalone input fragment. Expects amsmath,amssymb,amsthm, proposition,
% lemma,theorem environments. All macros below use standard LaTeX.
\section{Target mass and spectral compression}
\label{sec:allgamma-spectral}
This section extends the Gaussian spectral argument to unbounded
activations. Let $\phi(0)=0$ and let $\phi$ be $\lambda$-Lipschitz; replace
$\lambda$ by $\max\{1,\lambda\}$, so throughout $\lambda\ge1$.
For $\gamma\ge1$, the independent coordinate variances are
$N^{1-\gamma}/d$ in the first layer, $N^{-\gamma}$ internally, and
$N^{-1-\gamma}$ at the readout. Write $U_{N,g}$ for the globally
strict-margin event $g(x)z_w(x)>1/2$ for every cube input.
The normalized energy is
\[
 \kappa_N(w)=\frac dN\|W^1\|_F^2+
 \sum_{s=2}^{L-1}\|W^s\|_F^2+N\|v\|_2^2.
\]
No uniform bound on $\phi$ or on individual hidden activations is imposed.

\subsection{A target cylinder for linearly growing activations}
Let a depth-$L$, width-$C$ teacher have every weight bounded by $C$ and
strict output margin $1/2$. Multiplying only its readout by three gives
strict margin $3/2$. Set
\begin{align*}
 M&=(d+1)(\lambda C^2)^{L-1},&
 b&=2+\sqrt{2\log(2L)},\\
 B&=2+\lambda(\sqrt2 C^2+\sqrt C+b),&
 E&=\lambda M(1+LC)B^{L-2},\\
 \rho&=\left[4\{(3\sqrt2 C^2+1)E+M\}\right]^{-1},\\
 B_0&=d(d+1)C^2+2(L-2)C^4+36C^4,&
 D_0&=d+2+(L-2)C,\\
 A&=2B_0+4D_0/\rho+2.
\end{align*}
These quantities are independent of learner width and polynomially bounded
in $d,C$ for fixed $L,\lambda$. In particular $0<\rho\le1$ and $A\ge1$.

\begin{proposition}[Unbounded-activation target mass]
\label{prop:unbounded-target-mass}
At every $N\ge C$ and every $\gamma\ge1$, the raw Gaussian probability
of the target event satisfies
\[
 \mu_{N,\gamma}(U_{N,f})\ge \exp(-AN^\gamma/2).
\]
\end{proposition}
\begin{proof}
Partition the $N$ neurons into $C$ groups of sizes $k_a$ differing by at
most one. Then $N/k_a\le2C$ and $k_a/k_b\le2$. Clone the teacher by
copying each first-layer row throughout its group, replacing an internal
weight $B^s_{ba}$ by $B^s_{ba}/k_a$ on edges from group $a$ to group $b$,
and replacing readout weight $v_b$ by $v_b/k_b$. This preserves the exact
real-valued function for any activation. Every cloned hidden state is
group-constant. Let $P$ project onto that subspace.

Because $|\phi(u)|\le\lambda|u|$, all teacher hidden coordinates are
bounded in magnitude by $M$: the first layer is bounded by
$\lambda C(d+1)$, and each next hidden layer multiplies this bound by at
most $\lambda C^2$. The cloned internal operator norms are at most
$\sqrt2 C^2$, and the scaled readout norm is at most
$\sqrt N\|v_*\|_2\le3\sqrt2 C^2$.

Put $t=N^{(1-\gamma)/2}\le1$. In the first layer constrain each
$\sqrt d W^1_{ij}$ within $\rho t$ of its cloned value. In each internal
row, constrain its $C$ coordinates along the normalized group indicators,
after multiplying the row by $\sqrt N$, within $\rho t$ of their cloned
values. Finally constrain each $Nv_j$ within $\rho t$ of its cloned
value. All these are independent scalar Gaussian coordinates of variance
$t^2$. Write the resulting internal matrices as
\[
 W^s=W_*^s+R^s+E^s,\qquad R^s=R^sP,\quad E^sP=0.
\]
The constrained error obeys $\|R^s\|_F\le\rho t\sqrt C$; the free
part is $(t/\sqrt N)G^s(I-P)$, independently of the constrained
coordinates, with standard Gaussian $G^s$. The Gaussian operator-norm
bound and a union bound give conditional probability at least $1/2$ that
$\|E^s\|_{\rm op}\le bt$ for all internal layers.

On this event put
$e_s(x)=N^{-1/2}\|h_w^s(x)-h_*^s(x)\|_2$. Uniformly over the cube,
\[
 e_1\le\lambda\rho(d+1),\qquad
 e_s\le B e_{s-1}+\lambda\rho\sqrt C\,M.
\]
The second estimate uses $E^sh_*^{s-1}=0$, the bound on the teacher hidden
states, and Lipschitz continuity. It does not bound the perturbed states
by one. Since $M\ge d+1$ and $B\ge1$, induction gives
$e_{L-1}\le\rho E$. The readout error has scaled norm at most $\rho$,
so
\[
 \sup_x|z_w(x)-z_*(x)|
 \le (3\sqrt2 C^2+1)e_{L-1}+\rho M\le1/4.
\]
Thus every such draw lies in $U_{N,f}$ with room to spare.

There are $ND_0$ constrained coordinates and their squared centers sum
to at most $NB_0$. For $Z\sim N(0,t^2)$,
\[
 \mathbb P\{|Z-a|\le\rho t\}
 \ge\frac{2\rho}{\sqrt{2\pi}}e^{-a^2/t^2-\rho^2}.
\]
Multiplying these bounds, retaining the independent operator-norm event,
and using $N\le N^\gamma$ gives
\[
 -\log\mu_{N,\gamma}(U_{N,f})
 \le N^\gamma B_0+ND_0\log\frac{e\sqrt{2\pi}}{2\rho}+\log2
 \le AN^\gamma/2.
\]
Here $\log(e\sqrt{2\pi}/(2\rho))\le2/\rho$ for $0<\rho\le1$.
\end{proof}

\subsection{Spectral truncation with controlled hidden RMS norms}
\begin{lemma}[Gaussian spectral tail]
\label{lem:unbounded-spectral-tail}
If $W$ is an $N\times N$ matrix with independent $N(0,\sigma^2)$ entries,
$a\ge8\sigma\sqrt N$, and $W_{>a}$ retains singular values greater than
$a$, then
\[
 \mathbb P\{\|W_{>a}\|_F^2>K\}\le e^{-K/(8\sigma^2)}.
\]
\end{lemma}
\begin{proof}
The Gaussian operator-norm inequality
$\mathbb P\{\|W\|_{\rm op}>2\sigma\sqrt N+\sigma u\}\le e^{-u^2/2}$
shows that the operator-norm ball of radius $a/2$ has probability greater
than $1/2$. Its Frobenius distance from $W$ is
$D(W)=[\sum_j(s_j(W)-a/2)_+^2]^{1/2}$. Gaussian isoperimetry gives
$\mathbb P\{D(W)>r\}\le e^{-r^2/(2\sigma^2)}$.
Since $s\le2(s-a/2)$ when $s>a$, one has
$\|W_{>a}\|_F\le2D(W)$. Take $r=\sqrt K/2$.
\end{proof}

For $K\ge1$, define
\[
 a_K=\left[16L\lambda^L(K+1)(\sqrt K+1)^{L-2}\right]^{-1},
 \qquad J_K=(L+3)K.
\]
Use the auxiliary class $\mathcal F(J)$ from Corollary~\ref{cor:count};
it includes the Boolean functions admitting a strict margin-$1/2$,
depth-$L$, fixed-$\phi$ implementation of energy at most $J$ at some
finite width.

\begin{proposition}[Aggregate functional tail]
\label{prop:unbounded-aggregate}
Suppose $\gamma>1$ and
\[
 N^{\gamma-1}\ge8(d+1),\qquad
 N^{\gamma-1}\ge64a_K^{-2}.
\]
Then
\[
 \mu_{N,\gamma}\!\left(\bigcup_{g\notin\mathcal F(J_K)}U_{N,g}\right)
 \le L e^{-KN^\gamma/8}.
\]
\end{proposition}
\begin{proof}
Outside an event of probability at most $Le^{-KN^\gamma/8}$,
\[
 \frac dN\|W^1\|_F^2\le K,\quad N\|v\|_2^2\le K,
 \quad\|(W^s)_{>a_K}\|_F^2\le K\quad(2\le s<L).
\]
Indeed, the first and last random quantities are respectively
$N^{-\gamma}\chi^2_{N(d+1)}$ and $N^{-\gamma}\chi^2_N$.
For $u\ge8D$, Markov's inequality with exponential parameter $3/8$ gives
$\mathbb P\{\chi_D^2>u\}\le 2^D e^{-3u/8}\le e^{-u/4}$; the first width condition
therefore suffices. The second width condition permits
Lemma~\ref{lem:unbounded-spectral-tail} on each internal matrix.

Replace each internal matrix by its $>a_K$ part. Retained operator norms
are at most $\sqrt K$ and discarded operator norms at most $a_K\le1$.
For the original hidden states write
$M_s=N^{-1/2}\sup_x\|h_w^s(x)\|_2$. The zero-at-zero and Lipschitz
assumptions imply
\[
 M_s\le\sqrt{2K}\,\lambda^s(\sqrt K+1)^{s-1}.
\]
In particular individual neurons need not be bounded, but their normalized
Euclidean norm is controlled uniformly over inputs. If $e_s$ is the
normalized Euclidean difference between original and truncated hidden
states, then
\[
 e_1=0,\qquad e_s\le\lambda\sqrt K\,e_{s-1}+\lambda a_K M_{s-1}.
\]
Induction gives, for $s\ge2$,
\[
 e_s\le a_K(s-1)\sqrt{2K}\lambda^s(\sqrt K+1)^{s-2}.
\]
The output discrepancy is consequently at most
\[
 a_K(L-2)\sqrt2K\lambda^{L-1}(\sqrt K+1)^{L-3}<1/8
\]
when $L\ge3$. It is zero when $L=2$.
A globally margin-separated original draw therefore retains strict margin
$3/8$ after truncation. Doubling the readout gives strict margin $3/4$
and energy at most $K+(L-2)K+4K=J_K$. Its Boolean function belongs
to $\mathcal F(J_K)$, proving the claimed aggregate tail.
\end{proof}

The small spectral cutoff requires a sufficiently strong scale. At mean
field, the Gaussian bulk has operator norm of constant order; the next
sections control its functional effect through covariance reconstruction.
